% ============================================================================
%   CHAPTER 2 — STATE OF THE ART
%   Aivancity: 10-15 pages, thematic organization, ≥15 peer-reviewed refs
% ============================================================================

\chapter{State of the Art}
\label{chap:sota}

This chapter reviews the state of the art at the intersection of three 
research areas that jointly frame the present study: patient--clinical 
trial matching, probabilistic calibration of classifiers, and the 
evaluation of medically specialized large language models. Rather than 
surveying each area in isolation, we organize the review around the 
specific gap that motivates this thesis --- the absence of systematic 
calibration audits in LLM-based clinical trial matching --- and use 
each section to build the position of this work with respect to the 
closest prior contributions. Section~\ref{sec:sota-positioning} closes 
the chapter with a comparative table and an explicit statement of the 
contribution.

% ----------------------------------------------------------------------------
\section{Patient--Clinical Trial Matching}
\label{sec:sota-matching}

\subsection{Clinical Context and Recruitment Bottlenecks}
\label{subsec:sota-clinical}

Clinical trial recruitment is one of the most persistent bottlenecks in 
translational medicine. Despite widespread interest from patients and 
physicians, less than 5\% of adult cancer patients enroll in a clinical 
trial, and 80\% of trials fail to meet their recruitment 
targets~\cite{jin2024trialgpt}. The mismatch between patients and 
appropriate trials is caused, in part, by the sheer volume of open 
studies (over 400{,}000 registered on ClinicalTrials.gov at the time of 
writing) and the complexity of their eligibility criteria, which 
combine numeric thresholds, medical history requirements, and drug 
exclusions in dozens of criteria per trial. Manual screening of a 
patient against this space of criteria is labor-intensive and 
error-prone, and automating even a fraction of this process would 
substantially accelerate recruitment.

Automated matching systems have thus been a topic of active research 
for over a decade. The field has passed through three distinct waves, 
each corresponding to a different technological substrate.

\subsection{Pre-LLM Approaches: Rule-Based and Embedding Systems}
\label{subsec:sota-prelm}

The first wave of matching systems relied on structured information 
extraction. EliXR~\cite{weng2011elixr}, EliIE~\cite{kang2017eliie}, and 
Criteria2Query~\cite{yuan2019criteria2query} transformed free-text 
eligibility criteria into structured OMOP or SQL queries, matched 
against structured electronic health records. These systems produce 
deterministic Boolean outputs (match or no match) and provide no 
probability estimates or uncertainty quantification. They perform well 
on criteria that map cleanly to structured concepts (age, diagnoses, 
laboratory values) but struggle with criteria requiring nuanced 
interpretation.

The second wave introduced end-to-end neural matching. 
DeepEnroll~\cite{zhang2020deepenroll} encoded criteria with BERT and 
patient EHRs with hierarchical embeddings, framing matching as textual 
entailment and reporting PR-AUC and F1 scores. 
COMPOSE~\cite{gao2020compose} added a cross-modal pseudo-siamese 
network with a composite inclusion/exclusion loss, achieving 98.0\% 
patient--criterion AUC and 83.7\% patient--trial accuracy. Neither 
system reported calibration metrics. This omission is not accidental: 
the primary evaluation targets of these systems were discrimination 
metrics (AUC, F1) rather than the reliability of the probabilistic 
outputs.

\subsection{The LLM Wave: 2023--2026}
\label{subsec:sota-llmwave}

The third and current wave dates from the release of GPT-4 and the 
subsequent proliferation of open-source large language models. The 
most influential system in this wave is 
TrialGPT~\cite{jin2024trialgpt}, published in \emph{Nature 
Communications} in 2024. TrialGPT is a three-module pipeline: a 
retrieval stage recovers more than 90\% of relevant trials in less 
than 6\% of the collection using a hybrid BM25 and MedCPT retriever; a 
matching stage emits a per-criterion ordinal label 
(``included / not included / no relevant information''); and a ranking 
stage aggregates criterion labels into trial-level scores. TrialGPT 
reports a criterion-level accuracy of 87.3\% and correlation patterns 
against the ground truth, but does not report any calibration 
metric --- no ECE, no Brier score, no reliability diagram, and no 
monotonicity check across the ordinal scale.

TrialGPT is not an isolated case. RECTIFIER~\cite{unlu2024rectifier}, 
published in \emph{NEJM AI}, applies a retrieval-augmented GPT-4 to 
the COPILOT-HF heart-failure trial and reports 97.9--100\% accuracy 
with Matthews correlation coefficients between 0.837 and 1.00. The 
system has been deployed clinically at Mass General Brigham, yet its 
outputs are reported only as accuracy and MCC, without probabilistic 
reliability assessment. PRISM and its variant 
OncoLLM~\cite{gupta2024oncollm}, published in \emph{npj Digital 
Medicine}, introduce a weighted-tier scoring method with concept-level 
importance and report top-$k$ accuracy and nDCG (68\% vs.\ 62.6\% for 
GPT-3.5) but again omit calibration. Trial-LLaMA distills open-source 
Llama on synthetic criterion labels~\cite{nievas2024triallama} and 
reports accuracy alone. Zero-shot matching approaches by Wornow et 
al.~\cite{wornow2024zeroshot} achieve 0.93 micro-F1 on the n2c2 2018 
dataset and elicit a self-declared confidence per criterion --- the 
closest any published system comes to uncertainty quantification --- 
but the self-declared confidence is never calibrated against ground 
truth. More recent systems including 
TrialMatchAI~\cite{abdallah2026trialmatchai}, multimodal pipelines by 
Callies et al.~\cite{callies2025multimodal}, and cross-modal trial-site 
selectors extend the scope of matching but continue the pattern of 
accuracy-only reporting.

The definitive scoping review of the field, published by Chen et 
al.~\cite{chen2025scoping} in \emph{JCO Clinical Cancer Informatics}, 
screened 2{,}357 papers and included 24 LLM-based matching systems, 
of which 21 were published in 2024. The review confirms that 
\emph{none} of these systems reports calibration metrics: no ECE, no 
Brier score, no reliability diagram, no calibration-in-the-large, and 
no slope. This gap defines the primary methodological opportunity 
addressed by the present thesis.

% ----------------------------------------------------------------------------
\section{Probabilistic Calibration of Classifiers}
\label{sec:sota-calibration}

\subsection{Foundations and the Modern Miscalibration Problem}
\label{subsec:sota-foundations}

Probability calibration --- the requirement that predicted 
probabilities match empirical accuracies --- has a long history in 
statistical prediction. Historical foundations include Platt 
scaling~\cite{platt1999} for logistic regression outputs, isotonic 
regression~\cite{zadrozny2002isotonic} for non-parametric monotone 
calibration, and the systematic comparisons of Niculescu-Mizil and 
Caruana~\cite{niculescu2005}. For decades, calibration was treated as 
a solved problem: shallow classifiers such as logistic regression, 
naïve Bayes, and calibrated support vector machines were well 
calibrated after appropriate post-hoc correction.

This picture changed with the rise of deep learning. Guo, Pleiss, Sun, 
and Weinberger~\cite{guo2017calibration} showed in 2017 that modern 
deep neural networks, despite their superior accuracy, are 
\emph{systematically overconfident} relative to shallower architectures: 
a network with 85\% accuracy may assign 99\% confidence to its 
predictions, including its errors. This paper introduced the Expected 
Calibration Error (ECE) as the standard metric for measuring 
miscalibration in modern classifiers, and demonstrated the 
effectiveness of temperature scaling, a single-parameter post-hoc 
method that does not modify argmax predictions.

Since 2017, the calibration literature has expanded rapidly. 
Alternative metrics have been proposed to address specific limitations 
of ECE: adaptive-ECE, class-wise ECE, and the more refined estimators 
of Nixon et al.~\cite{nixon2019measuring} and Kumar, Liang, and 
Ma~\cite{kumar2019verified}. The Brier score~\cite{brier1950} provides 
a joint measure of calibration and sharpness. Beyond ECE, 
distribution-free confidence intervals for calibration have been 
studied by Gupta and Ramdas~\cite{gupta2021distributionfree}.

\subsection{Clinical Importance of Calibration}
\label{subsec:sota-clinical-calib}

The clinical importance of calibration has been argued particularly 
forcefully by Van Calster et al.~\cite{vancalster2019calibration}, 
whose 2019 \emph{BMC Medicine} article titled ``Calibration: the 
Achilles heel of predictive analytics'' establishes calibration as a 
prerequisite for clinical deployment of any predictive model. Van 
Calster and colleagues distinguish four progressive levels of 
calibration --- calibration in the large, weak calibration, moderate 
calibration, and strong calibration --- and argue that most published 
clinical prediction models fail even the weakest level. They also 
introduce the threshold ECE $< 0.05$ as a practical benchmark for 
clinically usable calibration, which the present thesis adopts as the 
rejection criterion for Hypothesis H1 (Section~\ref{sec:met-hypotheses}).

The 2024 update of the TRIPOD guidelines, known as 
TRIPOD+AI~\cite{collins2024tripod}, formalizes calibration as a 
required reporting element for any AI-based clinical prediction model. 
Under TRIPOD+AI, a paper introducing a clinical AI model must report 
calibration curves, calibration-in-the-large, and calibration slopes; 
failure to do so is now grounds for methodological rejection by 
several leading medical journals. This regulatory shift makes the 
absence of calibration reporting in LLM-based trial matching systems 
increasingly untenable.

\subsection{Recalibration Methods}
\label{subsec:sota-recalibration}

Post-hoc recalibration methods apply a monotone transformation to 
uncorrected confidences after the classifier has been trained. The 
three methods evaluated in the present thesis --- temperature scaling, 
isotonic regression, and Platt scaling --- represent the classical 
choices; more elaborate methods include beta 
calibration~\cite{kull2017beta}, Dirichlet 
calibration~\cite{kull2019dirichlet}, and the matrix/vector scaling 
family. Intrinsic methods, which modify the training objective itself, 
include label smoothing~\cite{muller2019labelsmoothing} and focal 
loss~\cite{mukhoti2020focal}. Distribution-free calibration approaches 
combine post-hoc scaling with conformal 
guarantees~\cite{gupta2021distributionfree}.

Complementary to point-estimate calibration, split conformal 
prediction~\cite{vovk2005, angelopoulos2023conformal} produces 
prediction sets with a guaranteed marginal coverage of $1 - \alpha$ 
without any distributional assumption on the underlying model. The 
guarantees hold under exchangeability of calibration and test data, 
which is satisfied in the present study by the split protocol 
described in Section~\ref{subsec:met-data-split}. Selective 
classification --- the option to abstain from prediction --- has been 
formalized by Geifman and 
El-Yaniv~\cite{geifman2017selective, geifman2019selectivenet} through 
risk--coverage curves, which are used in Section~\ref{sec:res-h3h4} to 
test Hypothesis H4.